% ch04_u06 —— 原书 p.172–p.181（4.5 全部：仿射 Grassmann 流形）
% 承接 p.171 末的命题陈述之 (ii)（本单元不写 \bdchap）。
\SourcePage{172}

(ii) $\operatorname{Aut}_Z O$ 与 $\widetilde{G}(K)$ 在 $\lambda_{\mathcal{L}}^{\wedge}$ 上的作用，定义了 $\operatorname{Aut}_Z O\rtimes\widetilde{G}(K)$ 的一个作用。

\emph{证明.} 把 $\mathcal{L}\in Z\operatorname{tors}_{\theta}(X)$ 写成 $\mathcal{L}=\mathcal{E}\cdot\mathcal{L}_0$，其中 $\mathcal{E}\in Z\operatorname{tors}(X)$，$\mathcal{L}_0\in\omega^{1/2}(X)$。按定义，$\lambda_{\mathcal{L}}=l_{\mathcal{E}}\otimes\lambda'_{\mathcal{L}_0}$（$\operatorname{Bun}_G$ 上 $\mu_\infty$-torsor $l_{\mathcal{E}}$ 的定义见 4.1.4 或 4.1.6）。

把 $\mathcal{L}^{\mathrm{loc}}$ 看作 $\omega^{1/2}(O)$ 的对象（这是可以的：$Z\operatorname{tors}_{\theta}(O)$ 与 $\omega^{1/2}(O)$ 都恰有一个同构类）。用 $\mathcal{L}_0$ 与 $\mathcal{L}^{\mathrm{loc}}$ 构造 $X^{\wedge}_2$、$M^{\wedge}_2$ 与 $\lambda^{\wedge}_{\mathcal{L}_0}$（见 4.4.14）。

把 $\mathcal{E}$ 看作 $Z$-覆叠 $\mathcal{E}\to X$。置 $X^{\wedge}_{\mathcal{E}}:=\mathcal{E}\times_X X^{\wedge}$，$M^{\wedge}_{\mathcal{E}}:=\mathcal{E}\times_X M^{\wedge}$，其中 $X^{\wedge}$ 与 $M^{\wedge}$ 的含义同 2.6.5 与 2.8.3。记 $l^{\wedge}_{\mathcal{E}}$ 为 $l_{\mathcal{E}}$ 到 $M^{\wedge}_{\mathcal{E}}$ 的拉回。

置 $M^{\wedge}_{\mathcal{E},2}:=\mathcal{E}\times_X M^{\wedge}_2$。有平展覆叠 $M^{\wedge}_{\mathcal{E},2}\to M^{\wedge}_2$、$M^{\wedge}_{\mathcal{E},2}\to M^{\wedge}_{\mathcal{E}}$，以及 $p:M^{\wedge}_{\mathcal{E},2}\to M^{\wedge}_Z$。显然 $p^*\lambda^{\wedge}_{\mathcal{L}}$ 是 $l^{\wedge}_{\mathcal{E}}$ 与 $\lambda^{\wedge}_{\mathcal{L}_0}$ 到 $M^{\wedge}_{\mathcal{E},2}$ 的拉回的张量积。下面分别考察 $l^{\wedge}_{\mathcal{E}}$ 与 $\lambda^{\wedge}_{\mathcal{L}_0}$。

半直积 $\operatorname{Aut} O\rtimes G(K)$ 作用在 $M^{\wedge}_{\mathcal{E}}$ 上；$\operatorname{Aut} O$ 在 $M^{\wedge}_{\mathcal{E}}$ 上的作用，典范地提升为它在 $l^{\wedge}_{\mathcal{E}}$ 上的作用（参看 4.4.14 或 2.8.5）。$G(K)$ 的作用对象是 $l^{\wedge}_{\mathcal{E}}$ 于 $X^{\wedge}_{\mathcal{E}}$ 各点纤维上的限制（见 4.1.7）。容易看出，这些作用来自 $\operatorname{Aut} O\rtimes G(K)$ 在 $l^{\wedge}_{\mathcal{E}}$ 上的一个作用。

另一方面，由 4.4.14，有 $\operatorname{Aut}_2 O\rtimes\widetilde{G}(K)$ 在 $\lambda^{\wedge}_{\mathcal{L}_0}$ 上的典范作用。

于是得到 $\operatorname{Aut}_2 O\rtimes\widetilde{G}(K)$ 在 $p^*\lambda^{\wedge}_{\mathcal{L}}$ 上的作用。它与 $\operatorname{Aut}_2 O$ 在 $\lambda^{\wedge}_{\mathcal{L}}$ 上的作用相容，也与 $\widetilde{G}(K)$ 在 $\lambda^{\wedge}_{\mathcal{L},\hat x}$（$\hat x\in X^{\wedge}_Z$）上的作用相容。

$p$ 是平展满射，故 $\operatorname{Aut}_2 O\rtimes\widetilde{G}(K)$ 在 $p^*\lambda^{\wedge}_{\mathcal{L}}$ 上的作用下降为它在 $\lambda^{\wedge}_{\mathcal{L}}$ 上的作用。$\operatorname{Aut}_Z O$ 由 $\operatorname{Aut}_2 O$ 与 $Z$ 生成，故只需证 $Z\subset\operatorname{Aut}_Z O$ 在 $\lambda^{\wedge}_{\mathcal{L}}$ 上的作用与 $\widetilde{G}(K)$ 的作用相容。这一点显然：对每个 $\hat x\in X^{\wedge}_Z$，$Z$ 与 $\widetilde{G}(K)$ 在 $\lambda^{\wedge}_{\mathcal{L},\hat x}$ 上的作用都相容。\hfill$\Box$

\bdsec{4.5. 仿射 Grassmann 流形}

仿射 Grassmann 流形 $\mathcal{GR}$ 是 fpqc-商 $G(K)/G(O)$，其中 $O=\mathbb{C}[[t]]$，$K=\mathbb{C}((t))$。\SourcePage{173}本节回顾 $\mathcal{GR}$ 的一些基本性质。4.6 节将构造并研究\emph{局部 Pfaffian 丛}；它是 $\mathcal{GR}$ 上的线丛。

仿射 Grassmann 流形在主要定理 5.2.6 的证明中起关键作用。不过读者暂时可以跳过本节。

在 4.5.1 – ? 中，$G$ 表示任一连通仿射代数群。连通性是无害的假设，因为 $G(K)/G(O)=G^0(K)/G^0(O)$，其中 $G^0$ 是 $G$ 的连通分支。
% [Tutor Note: 原书 p.173「In 4.5.1 – ?」中的交叉引用只印出一个问号、无编号，底本如此，保留原样。]

\bdp{4.5.1} \emph{定理.}
\begin{enumerate}[label=(\roman*),leftmargin=2.2em]
\item fpqc-商 $G(K)/G(O)$ 是归纳有限型的归纳概形。
\item $G(K)/G(O)$ 形式光滑。\footnote{形式光滑的定义见 7.11.1。}
\item 投影 $p:G(K)\to G(K)/G(O)$ 在 Zariski 拓扑下局部有截面。
\item $G(K)/G(O)$ 是归纳真的，当且仅当 $G$ 约化。
\item $G(K)$（等价地，$G(K)/G(O)$）既约，当且仅当 $\operatorname{Hom}(G,\mathbb{G}_m)=0$。
\end{enumerate}

\emph{注.} 这个定理是众所周知的。下面给出的证明，其实质部分只是引用 Faltings、Beauville、Laszlo 与 Sorger 的工作。

\emph{证明.} 对 $G=GL_n$，(i) 与 (iv) 成立。事实上，存在归纳真的归纳概形 $Gr(K^n)$，它参数化 $K^n$ 中的 c-格（细节见 7.11.2(iii)）。$GL_n(K)/GL_n(O)$ 等同于 $Gr(K^n)$ 中参数化 $O$-不变 c-格的闭子归纳概形。为对任意 $G$ 证明 (i) 与 (iv)，需要下面的引理。

\emph{引理.} 设 $G_1\subset G_2$ 是仿射代数群，且商 $U:=G_1\backslash G_2$ 拟仿射，即 $U$ 是仿射概形\SourcePage{174}$Z$ 的开子概形。设 fpqc-商 $G_2(K)/G_2(O)$ 是归纳有限型的归纳概形。则 $G_1(K)/G_1(O)$ 亦然，且态射
\begin{equation}\tag{221}
G_1(K)/G_1(O)\to G_2(K)/G_2(O)
\end{equation}
是局部闭嵌入。若 $U$ 仿射，则 (221) 是闭嵌入。

读者可用 4.5.2 中 $G(K)/G(O)$ 的整体解释轻易证明引理。我们更愿意给出局部证明。

\emph{证明.} 考虑态射 $f:G_1(K)\to Z(K)$。显然 $Z(O)$ 是 $Z(K)$ 的闭子概形，$U(O)$ 是 $Z(O)$ 的开子概形。故 $Y:=f^{-1}(U(O))$ 是 $G_2(K)$ 的局部闭子归纳概形；当 $U$ 仿射时它是闭的。显然 $Y\cdot G_2(O)=Y$，故 $Y$ 是某个局部闭子归纳概形 $Y'\subset G_2(K)/G_2(O)$ 的原像；当 $U$ 仿射时 $Y'$ 是闭的。由 $G_1(K)\subset Y$，有自然态射
\begin{equation}\tag{222}
G_1(K)\to Y'.
\end{equation}

我们断言 (222) 是 $G_1(O)$-torsor（$G_1(O)$ 经右平移作用在 $G_1(K)$ 上），从而 $G_1(K)/G_1(O)=Y'$。为看出 (222) 是 $G_1(O)$-torsor，注意态射 $Y\to Y'$ 是 $G_2(O)$-torsor，态射 $\varphi:Y\to U(O)=G_1(O)\backslash G_2(O)$ 是 $G_2(O)$-等变的，且 $G_1(K)=\varphi^{-1}(\bar e)$，其中 $\bar e\in G_1(O)\backslash G_2(O)$ 是 $e\in G_2(O)$ 的像。\hfill$\Box$

下面证明对任意 $G$ 的 (i) 与 (iv)。取嵌入 $G\hookrightarrow GL_n$。若 $G$ 约化，则 $GL_n/G$ 仿射，故由引理知 $G(K)/G(O)$ 是归纳真的归纳概形。对任意 $G$，我们将构造嵌入 $i:G\hookrightarrow G':=GL_n\times\mathbb{G}_m$，使 $G'/i(G)$ 拟仿射；由此推出 (i)。为构造 $i$，取 $GL_n$-模 $V$，使 $G\subset GL_n$ 是某个一维子空间 $l\subset V$ 的稳定化子。$G$ 在 $l$ 中的作用由某个 $\chi:G\to\mathbb{G}_m$ 给出。定义 $i:G\hookrightarrow G':=GL_n\times\mathbb{G}_m$，$i(g)=(g,\chi(g)^{-1})$。为证 $G'/i(G)$ 拟仿射，把 $V$ 看作 $G'$-模（$\lambda\in\mathbb{G}_m$ 作用为乘 $\lambda$），并注意非零 $v\in l$ 在 $G'$ 中的稳定化子等于 $i(G)$。故 $G'/i(G)\simeq G'v$，而 $G'v$ 拟仿射。

\SourcePage{175}

下面完成 (iv) 的证明。若 $G(K)/G(O)$ 是归纳真的，且 $G'$ 是 $G$ 的正规子群，则按引理 $G'(K)/G'(O)$ 也是归纳真的。显然 $\mathbb{G}_a(K)/\mathbb{G}_a(O)$ 不是归纳真的。因此只有当 $G$ 约化时，$G(K)/G(O)$ 才归纳真。

为证 (iii)，只需证 $p:G(K)\to G(K)/G(O)$ 在任一 $\mathbb{C}$-点 $x\in G(K)/G(O)$ 的邻域上有截面（这里用到由 (i) 知 $\mathbb{C}$-点在 $G(K)/G(O)$ 中稠密）。由 $p$ 的 $G(K)$-等变性，可化到 $x$ 是 $e\in G(K)$ 的像的情形。于是须构造含 $e$ 的子归纳概形 $\Gamma\subset G(K)$，使态射
\begin{equation}\tag{223}
\Gamma\times G(O)\to G(K),\qquad (\gamma,g)\mapsto\gamma g
\end{equation}
是开浸入。按 Faltings [Fal94, p.350–351]，当 $\Gamma$ 的 $R$-点集由
\[
\Gamma(R)=\operatorname{Ker}\bigl(G(R[t^{-1}])\xrightarrow{f}G(R)\bigr)\subset G(R((t)))=G(R\hat\otimes K)
\]
定义时，(223) 是开浸入，其中 $f$ 是在 $t=\infty$ 处的求值。这个断言的证明属于 Beauville 与 Laszlo（[BLa94] 中的命题 1.11）。它以 $X=\mathbb{P}^1$ 情形下 $G(K)/G(O)$ 的整体解释为基础（见 4.5.2），也以 $\mathbb{P}^1$ 上 $G$-丛的下述性质为基础：对 $S\times\mathbb{P}^1$ 上的 $G$-丛 $\mathcal{F}$，使 $\mathcal{F}$ 在 $s\times\mathbb{P}^1$ 上的限制平凡的那些点 $s\in S$ 构成 $S$ 的\emph{开}子集（事实上，$H^1(\mathbb{P}^1,\mathcal{O}\otimes\mathfrak{g})=0$，$\mathfrak{g}:=\operatorname{Lie} G$）。

由 (iii) 推出\footnote{事实上，不用 (iii) 也能证明 (ii)。} (ii)。$G(K)$ 形式光滑，故由 (iii) 知 $G(K)/G(O)$ 的每一点都有形式光滑的邻域。又 $G(K)/G(O)$ 是归纳有限型的，故得 (ii)。

还需考虑 (v)。$G(O)$ 是既约的。故 $G(K)$ 既约当且仅当 $G(K)/G(O)$ 既约。Laszlo 与 Sorger 证明：若 $\operatorname{Hom}(G,\mathbb{G}_m)=0$，则 $G(K)/G(O)$ 既约（见 [La-So] 中命题 4.6 的证明）；其证明基于\SourcePage{176}Shafarevich 的一个定理。若 $\operatorname{Hom}(G,\mathbb{G}_m)\ne 0$，则存在态射 $f:\mathbb{G}_m\to G$ 与 $\chi:G\to\mathbb{G}_m$，使 $\chi f=\varphi_n$，$n\ne 0$，其中 $\varphi_n(\lambda):=\lambda^n$。由 $\varphi_n$ 诱导的态射 $\mathbb{G}_m(K)\to\mathbb{G}_m(K)$ 的像不含于 $\mathbb{G}_m(K)_{red}$ 中，故 $G(K)$ 不是既约的。\hfill$\Box$

\bdp{4.5.2} 设 $X$ 是 $\mathbb{C}$ 上的连通光滑射影曲线，$x\in X(\mathbb{C})$，$O_x$ 是 $x$ 处的完备化局部环，$K_x$ 是其分式域。按 Beauville–Laszlo（见 2.3.4），fpqc-商 $G(K_x)/G(O_x)$ 可解释为对 $(\mathcal{F},\gamma)$ 的模空间。这里 $\mathcal{F}$ 是 $X$ 上的主 $G$-丛，$\gamma:X\setminus\{x\}\to\mathcal{F}$ 是它的截面（即平凡化）。把 $(\mathcal{F},\gamma)$ 对应到 $\gamma/\gamma_x$ 在 $G(K_x)/G(O_x)$ 中的像；其中 $\gamma_x$ 是 $\mathcal{F}$ 在 $\operatorname{Spec} O_x$ 上的截面，$\gamma/\gamma_x$ 表示满足 $\gamma=g\gamma_x$ 的元素 $g\in G(K_x)$。我们已在 $\mathbb{C}$-点层面把 $G(K_x)/G(O_x)$ 与对 $(\mathcal{F},\gamma)$ 的模空间等同；对任意 $\mathbb{C}$-代数 $R$，读者容易对 $R$-点完成同样的事。

\bdp{4.5.3} 下面回顾 $G$ 的拓扑基本群的代数定义。记 $\pi_1^{\text{ét}}(G)$ 为 $G$ 在 Grothendieck 意义下的基本群。特征 $f:G\to\mathbb{G}_m$ 诱导态射 $\pi_1^{\text{ét}}(G)\to\pi_1^{\text{ét}}(\mathbb{G}_m)=\hat{\mathbb{Z}}(1)$，从而诱导态射 $f_*:(\pi_1^{\text{ét}}(G))(-1)\to\hat{\mathbb{Z}}$。记 $\pi_1(G)$ 为 $(\pi_1^{\text{ét}}(G))(-1)$ 中满足下述条件的 $\alpha$ 的集合：对一切 $f\in\operatorname{Hom}(G,\mathbb{G}_m)$ 有 $f_*(\alpha)\in\hat{\mathbb{Z}}$。我们把 $\pi_1(G)$ 看作离散群。事实上，把 $G$ 换成它的极大约化商，$\pi_1(G)$ 不变。对约化 $G$，可把 $\pi_1(G)$ 等同于 $G$ 的余权群模去余根格所得的商。

对任意有限覆叠 $p:\tilde{G}\to G$，有上边缘映射 $G(K)\to H^1(K,A)=A(-1)$，$A:=\operatorname{Ker} p$。这些映射给出同态 $G(K)\to(\pi_1^{\text{ét}}(G))(-1)$。其像含于 $\pi_1(G)$ 中。于是我们构造了典范同态
\begin{equation}\tag{224}
\varphi:G(K)\to\pi_1(G),
\end{equation}
\SourcePage{177}其中 $G(K)$ 按朴素意义理解（即看作 $G$ 的 $K$-点群，或归纳概形 $G(K)$ 的 $\mathbb{C}$-点群）。(224) 在 $G(O)$ 上的限制平凡，故 (224) 诱导映射
\begin{equation}\tag{225}
G(K)/G(O)\to\pi_1(G),
\end{equation}
其中 $G(K)/G(O)$ 也按朴素意义理解。

现在把 $G(K)$ 与 $G(K)/G(O)$ 看作归纳概形。$G(K)/G(O)$ 的 $\mathbb{C}$-点集在 $G(K)/G(O)$ 中稠密；$G(K)$ 亦然。

\bdp{4.5.4} \emph{命题.}
\begin{enumerate}[label=(\roman*),leftmargin=2.2em]
\item 映射 (224) 与 (225) 局部常值。
\item 相应的映射
\begin{equation}\tag{226}
\pi_0(G(K))\to\pi_1(G),
\end{equation}
\begin{equation}\tag{227}
\pi_0(G(K)/G(O))\to\pi_1(G)
\end{equation}
是双射。
\end{enumerate}

\emph{证明.} 我们已用整体论证证明了 (i)（见 4.1.7 末的注）。同样的论证可改用 4.5.2 中 $G(K_x)/G(O_x)$ 的解释来表述：映射 (225) 等于自然映射 $G(K_x)/G(O_x)\to\operatorname{Bun}_G$ 与「第一陈类」映射 $c:\pi_0(\operatorname{Bun}_G)\to\pi_1(G)$ 的复合取负。(i) 的局部证明见 4.5.5。

现在证 (ii)。映射 $\pi_0(G(K))\to\pi_0(G(K)/G(O))$ 是双射（因为 $G$ 连通）。故只需考虑 (226)。因 $G$ 可表为约化群与幂幺群的半直积，可设 $G$ 约化。固定嘉当子群 $H\subset G$。有 $\pi_0(H(K))=\pi_1(H)$；复合 $\pi_0(H(K))\to\pi_0(G(K))\to\pi_1(G)$ 就是自然映射 $\pi_1(H)\to\pi_1(G)$，它是满射。故 (226) 也是满射。映射 $\pi_0(H(K))\to\pi_0(G(K))$ 是满射（对抽象群 $G(K)$ 用 Bruhat 分解）。因此，为证 (226) 单射，只需证自然态射\SourcePage{178}$f:\pi_0(H(K))\to\pi_1(G)$ 的核含于 $\pi_0(H(K))\to\pi_0(G(K))$ 的核中。因 $\operatorname{Ker} f$ 是余根格，只需证：对任意余根 $\gamma:\mathbb{G}_m\to H$，$\mathbb{G}_m(K)$ 在 $G(K)$ 中的像属于 $e\in G(K)$ 的连通分支。余根 $\mathbb{G}_m\to H$ 可延拓为态射 $SL(2)\to G$，故只需注意 $SL(2,K)$ 连通（因为 $SL(2,K)$ 中任一矩阵都可表为幂幺矩阵的乘积）。\hfill$\Box$

下一小节给出 4.5.4(i) 的局部证明。

\bdp{4.5.5} \emph{引理.} 设 $M=\operatorname{Spec} R$ 是连通仿射簇，$A$ 是有限交换群，$\alpha\in H^1_{\text{ét}}(\operatorname{Spec} R((t)),A)$。对 $x\in M(\mathbb{C})$，记 $\alpha(x)$ 为 $\alpha$ 在 $x$ 上方纤维 $\operatorname{Spec} R((t))\to\operatorname{Spec} R$ 上的限制，于是 $\alpha(x)\in H^1_{\text{ét}}(\operatorname{Spec}\mathbb{C}((t)),A)=A(-1)$。则 $\alpha(x)\in A(-1)$ 不依赖于 $x$。

\emph{证明.} 只需证：对任意光滑连通 $M'$ 与任意态射 $M'\to M$，$\alpha$ 到 $M'(\mathbb{C})$ 的拉回是常值\footnote{事实上，只须考虑那些光滑曲线 $M'$。}。于是可设 $M$ 光滑。置 $V:=\operatorname{Spec} R[[t]]$，$V':=\operatorname{Spec} R((t))$。可设 $A=\mu_n$。则 $\alpha$ 对应 $V'$ 上的 $\mu_n$-torsor，即 $V'$ 上的线丛 $\mathcal{A}$ 连同一个同构 $\psi:\mathcal{A}^{\otimes n}\xrightarrow{\ \sim\ }\mathcal{O}_{V'}$。因 $V$ 正则，$\mathcal{A}$ 延拓为 $V$ 上的线丛 $\tilde{\mathcal{A}}$。此时 $\psi$ 诱导同构 $\tilde{\mathcal{A}}^{\otimes n}\xrightarrow{\ \sim\ }t^k\mathcal{O}_V$，$k\in\mathbb{Z}$。显然 $\alpha(x)\in\mathbb{Z}/n\mathbb{Z}$ 就是 $k$ 的像。\hfill$\Box$

下面给出 4.5.4(i) 的局部证明。因 $G(K)/G(O)$ 是归纳有限型的，只需证：对每个连通仿射簇 $M=\operatorname{Spec} R$ 与任意态射 $f:M\to G(K)$，复合 $M(\mathbb{C})\to G(K)\to\pi_1(G)$ 是常值。对任意有限交换群 $A$，正合列 $0\to A\to\tilde{G}\to G\to 0$ 定义了映射 $\pi_1(G)\to A(-1)$；故只需证复合 $M(\mathbb{C})\to G(K)\to\pi_1(G)\to A(-1)$ 常值。为此把引理用于 $\alpha=\varphi^*\beta$，其中 $\varphi:\operatorname{Spec} R((t))\to G$ 对应于\SourcePage{179} $f:\operatorname{Spec} R\to G(K)$，而 $\beta\in H^1_{\text{ét}}(G,A)$ 是 $\tilde{G}$ 作为 $G$ 上 $A$-torsor 的类。

\emph{注.} 事实上可以证明：对 $\mathbb{C}$ 上任一仿射概形 $M=\operatorname{Spec} R$，所谓「Künneth 态射」
\begin{equation}\tag{228}
\begin{gathered}
H^1_{\text{ét}}(M,A)\oplus H^0(M,\mathbb{Z})\otimes H^1_{\text{ét}}(\operatorname{Spec}\mathbb{C}((t)),A)\to H^1_{\text{ét}}(M((t)),A),\\
M((t)):=\operatorname{Spec} R((t)),
\end{gathered}
\end{equation}
是同构（显然这蕴含引理）。对使 $A$ 的阶在 $R$ 中可逆的任意环 $R$，类似的陈述也成立。

\bdp{4.5.6} \emph{命题.} 设 $A\subset G$ 是有限中心子群，$G':=G/A$。
\begin{enumerate}[label=(\roman*),leftmargin=2.2em]
\item 态射 $G(K)/G(O)\to G'(K)/G'(O)$ 诱导 $G(K)/G(O)$ 与 $G'(K)/G'(O)$ 的若干连通分支之并之间的同构。
\item 态射 $G(K)\to G'(K)$ 是平展覆叠。
\end{enumerate}

\emph{注.} 由 4.5.4，(i) 中提到的分支由 $\operatorname{Im}(\pi_1(G)\to\pi_1(G'))$ 的元素标号。$G(K)$ 在 $G'(K)$ 中的像的连通分支也如此。

\emph{证明.} 显然 (i) 与 (ii) 等价。设 $G$ 半单（这等价于 $G'$ 半单）来证 (i)。此时态射 $f:G(K)/G(O)\to G'(K)/G'(O)$ 由 4.5.1(iv) 是归纳真的。由 4.5.4(i)，对满足 $f^{-1}(C)\ne\emptyset$ 的分支 $C\subset G'(K)/G'(O)$，$f$ 在其几何点\footnote{对 $\mathbb{C}$-点的陈述由 4.5.4(i) 立即得到。把 $\mathbb{C}$ 换成代数闭域 $E\supset\mathbb{C}$ 后 4.5.4 仍成立，故对 $E$-点的陈述也成立。}上的纤维恰含一点，且容易看出这些纤维既约。由 4.5.1(v)，$G'(K)/G'(O)$ 既约。故半单情形下 (i) 显然。

\SourcePage{180}

下面把 (ii) 的证明化到半单情形。可设 $A$ 循环。只需构造从 $G$ 到半单群 $G_1$ 的态射 $\rho$，使 $\rho|_A$ 单射且 $\rho(A)\subset G_1$ 是中心的（此时态射 $G(K)\to G'(K)$ 由 $G_1(K)\to G'_1(K)$ 经基变换得到，其中 $G'_1:=G_1/\rho(A)$）。构造 $G_1$ 与 $\rho$ 的办法如下。固定同构 $\chi:A\xrightarrow{\ \sim\ }\mu_n$。设 $V$ 是有限维 $G$-模，使 $Z$ 经 $\chi$ 作用在 $V$ 上。记 $W_{pq}$ 为 $p$ 份 $V$ 与 $q$ 份 $\operatorname{Sym}^{n-1}V^*$ 的直和。若 $p\cdot\dim V=q(n-1)\cdot\dim\operatorname{Sym}^{n-1}V$，则可置 $G_1:=SL(W_{pq})$（事实上，$GL(V)$ 在 $GL(W_{pq})$ 中的像含于 $SL(W_{pq})$。\hfill$\Box$
% [Tutor Note: 原书 p.180「(indeed, the image of … is contained in $SL(W_{pq})$.)」的左括号未闭合，底本如此，保留原样。]

\emph{注记.}
\begin{enumerate}[label=(\roman*),leftmargin=2.2em]
\item 命题 4.5.6 是 (228) 之双射性的直接推论。
\item 用 4.5.2 中 $G(K)/G(O)$ 的整体解释，容易证明命题 4.5.6。
\end{enumerate}

\bdp{4.5.7} 设 $G$ 约化。记 $G_{ad}$ 为 $G$ 模去其中心的商。置 $T:=G/[G,G]$，$G':=G_{ad}\times T$。则 $G'=G/A$，其中 $A\subset G$ 是某有限中心子群。故由 4.5.6，$G(K)/G(O)$ 可等同于 $G'(K)/G'(O)=G_{ad}(K)/G_{ad}(O)\times T(K)/T(O)$ 的某些连通分支之并。

$T(K)/T(O)$ 的结构相当简单。例如 $\mathbb{G}_m(K)/\mathbb{G}_m(O)$ 的既约部分是离散空间 $\mathbb{Z}$；$1\in\mathbb{G}_m(K)/\mathbb{G}_m(O)$ 的连通分支是以 $K/O$ 为李代数的形式群。

\bdp{4.5.8} 从现在起设 $G$ 约化，并置 $\mathcal{GR}:=G(K)/G(O)$。

回忆 $\mathcal{GR}$ 中的 $G(O)$-轨道由 $G$ 的支配余权标号；换言之，由 $P_+({}^LG):=$ ${}^LG$ 的支配权集标号。更确切地说，$\chi\in P_+({}^LG)$ 定义了态射 $\nu:\mathbb{G}_m\to G$ 的一个共轭类；按定义，$\operatorname{Orb}_\chi$ 是 $\nu(\pi)$ 在 $\mathcal{GR}$ 中的像的 $G(O)$-轨道，其中\SourcePage{181}$\pi$ 是 $O$ 的一个素元（这个像不依赖于 $\pi$ 的选取）。显然 $\operatorname{Orb}_\chi$ 不依赖于共轭类中 $\nu$ 的选取，故 $\operatorname{Orb}_\chi$ 有定义。按 [IM]，映射 $\chi\mapsto\operatorname{Orb}_\chi$ 是 $P_+({}^LG)$ 与 $\mathcal{GR}$ 中 $G(O)$-轨道集之间的双射。容易证明
\begin{equation}\tag{229}
\dim\operatorname{Orb}_\chi=(\chi,2\rho),
\end{equation}
其中 $2\rho$ 是 $G$ 的正根之和。

\emph{注.} 显然 $\operatorname{Orb}_\chi$ 是 $\operatorname{Aut}^0 O$-不变的。

\bdp{4.5.9} 已有 $\pi_0(\mathcal{GR})$ 与 $\pi_1(G)$ 之间的双射 (227)。设 $Z$ 是朗兰兹对偶群 ${}^LG$ 的中心。利用 $G$ 与 ${}^LG$ 的嘉当环面之间的对偶，把 $\pi_1(G)$ 等同于 $Z^\vee:=\operatorname{Hom}(Z,\mathbb{G}_m)$。于是 $\mathcal{GR}$ 的连通分支由 $Z^\vee$ 的元素标号。

\emph{注.} 含 $\operatorname{Orb}_\chi$ 的 $\mathcal{GR}$ 连通分支对应于 $\chi_Z\in Z^\vee$，其中 $\chi_Z$ 是 $\chi\in P^+({}^LG)$ 到 $Z$ 的限制。
% [Tutor Note: 原书 p.181 此处记作 $P^+({}^LG)$，而 p.180 同一对象记作 $P_+({}^LG)$，底本记号不一致，保留原样。]

\bdp{4.5.10} 存在典范态射 $\alpha:\mu_2\to Z$。当 $G$ 半单时，我们已用 (56) 定义过它。当 $G$ 约化时，这给出态射 $\mu_2\to Z'$，其中 $Z'$ 是 ${}^LG$ 的中心化子的中心；于是定义 $\alpha$ 为复合 $\mu_2\to Z'\hookrightarrow Z$。

按 4.4.4，对偶态射 $\alpha^\vee:\pi_1(G)\to\mathbb{Z}/2\mathbb{Z}$ 就是由伴随表示 $G\to SO(\mathfrak{g}_{ss})$ 给出的基本群之间的态射，其中 $\mathfrak{g}_{ss}:=[\mathfrak{g},\mathfrak{g}]$。

(227) 与 $\alpha^\vee$ 的复合定义了局部常值的\emph{奇偶函数}
\begin{equation}\tag{230}
p:\mathcal{GR}\to\mathbb{Z}/2\mathbb{Z}.
\end{equation}

若 (230) 把 $\mathcal{GR}$ 的连通分支映到 0（相应地 1），则称它是\emph{偶}的（相应地\emph{奇}的）。

\bdp{4.5.11} \emph{命题.} $\mathcal{GR}$ 的偶（相应地奇）连通分支上的一切 $G(O)$-轨道，维数都是偶（相应地奇）的。
